% ECONOMICS_PROBLEM: {"id": "E10-68", "title": "Sufficient heterogeneity in macro models", "jel_code": "E10", "source_id": "OP-0234"}
% !TeX program = xelatex
\documentclass[11pt,a4paper]{article}
\usepackage[margin=23mm]{geometry}
\usepackage{fontspec}
\setmainfont{Latin Modern Roman}
\usepackage{amsmath,amssymb,mathtools}
\usepackage{xcolor,enumitem,booktabs,longtable,array,xurl,hyperref}
\definecolor{muted}{HTML}{53636C}
\hypersetup{colorlinks=true,urlcolor=blue,linkcolor=blue}
\setlength{\parindent}{0pt}
\setlength{\parskip}{5pt}
\newcommand{\R}{\mathbb R}
\newcommand{\E}{\mathbb E}
\newcommand{\N}{\mathbb N}
\newcommand{\ind}{\mathbf 1}
\newcommand{\Var}{\operatorname{Var}}
\newcommand{\Cov}{\operatorname{Cov}}
\newcommand{\supp}{\operatorname{supp}}
\DeclareMathOperator*{\argmin}{arg\,min}
\DeclareMathOperator*{\argmax}{arg\,max}
\newcommand{\entrymeta}[3]{{\sffamily\small\color{muted}\textbf{\texttt{#1}}\quad|\quad #2\par #3\par}}
\newcommand{\Problemref}[1]{\texttt{#1}}
\newcommand{\JELref}[2]{\texttt{#2} (JEL \texttt{#1})}
\begin{document}
\section*{Sufficient heterogeneity in macro models}
\textbf{Problem ID:} \texttt{E10-68}\quad \textbf{JEL:} \texttt{E10}\par
% BEGIN ECONOMICS STATEMENT
\label{problem:OP-0234}
{\sffamily\small\color{muted}Original subject group: Business cycles and macro dynamics\par}
\label{definitions:problem:30007515}
\textbf{Audit classification:} Published research agenda. Heterogeneity matters to policy analysis; minimal continuous distributional compression is an editorial question. Verification concerns the cited version, not first priority or proof that this exact editorial specification remains unsolved on October 8, 2026.
\subsection*{Definitions and precise research target}
\paragraph{Editorial mathematical specification.} Let \(\mathcal M\) be a declared class of heterogeneous-agent economies sharing specified preferences, borrowing constraints, and production technologies, while their admissible asset, income, and firm-productivity distributions vary. For model \(M\in\mathcal M\), let \(\mathcal D_M\) be its admissible distribution set, \(\mu\) denote its cross-sectional distribution and \(R_M(\mu,a,h)\) its equilibrium response of output, consumption, employment, or inflation at horizon \(h\) to intervention \(a\). Interventions belong to a compact, explicitly bounded set \(\mathcal A\); horizons satisfy \(0\le h\le H\). Fix a response norm and accuracy \(\epsilon>0\).
Seek an economically interpretable statistic \(T(\mu)\in\mathbb R^d\), a transition law for that statistic, and a reduced equilibrium response \(\widehat R_M\) such that
\[
\sup_{M\in\mathcal M,\mu\in\mathcal D_M,a\in\mathcal A,\,0\le h\le H}
\|R_M(\mu,a,h)-\widehat R_M(T(\mu),a,h)\|\le\epsilon.
\]
Determine the smallest attainable dimension \(d\), or defensible upper and lower bounds, within a specified family of candidate statistics such as liquid-wealth bins, income exposure, and firm financing constraints. Explain which distributions and interventions invalidate a proposed compression. Establish approximation error separately from parameter-estimation error, and compare the compressed model with microeconomic evidence and held-out policy episodes. Without restrictions on \(\mathcal M\), finite-dimensional sufficiency need not exist; identifying those restrictions is part of this scoped target.

To prevent an arbitrary real number from encoding an entire distribution, put cross-sectional states in a fixed compact metric space \(X\), equip distributions with the Wasserstein distance \(W_1\), and restrict statistics to moments \(T_d(\mu)=(\int\varphi_j\,d\mu)_{j=1}^d\) from a fixed ordered family of bounded Lipschitz functions. Require reduced responses and their transition map to have a fixed Lipschitz bound independent of \(d\). Define \(d_\epsilon\) as the smallest positive integer for which the displayed error bound and specified transition-error bound are feasible, or \(+\infty\) if none exists. The benchmark response must use the same initial aggregate state, equilibrium selection, and intervention convention. Bounded horizons permit approximation even when exact Markov closure fails; exact sufficiency is a stronger property. For a prescribed tolerance \(\epsilon_T\ge0\), require \(\sup\|T_d(\mu_{t+1})-F_d(T_d(\mu_t),s_t,a_t)\|\le\epsilon_T\), with the supremum over the same admissible models, distributions, states, actions, and shocks for which the transition is defined.
\subsection*{Variants and assumption changes}
The following are \emph{editorial variants}, not additional published open conjectures.
\begin{enumerate}
\item \emph{Exact closure.} Set \(\epsilon=0\) and require \(T_d(\mu_{t+1})=F_d(T_d(\mu_t),s_t,a_t)\) for every admissible distribution and shock. Characterize primitives permitting this finite moment closure.
\item \emph{Local response compression.} Restrict \(\mu\) to a Wasserstein neighborhood of a stationary distribution and interventions to infinitesimal shocks. Seek compression of the derivative of the equilibrium response; this is weaker than global accuracy for large policy changes.
\end{enumerate}
\subsection*{References and source audit}
\begin{itemize}
\item Janet L. Yellen. \emph{Macroeconomic Research After the Crisis}. 2016. Locator: Section Heterogeneity. Primary text checked. \url{https://www.federalreserve.gov/newsevents/speech/yellen20161014a.htm}.
\item Adrien Auclert; Matthew Rognlie; Ludwig Straub. \emph{Fiscal and Monetary Policy with Heterogeneous Agents}. 2025. Locator: Section5, printed pp.19--21, especially Optimal policy; this does not state a minimal-compression conjecture. Primary text checked. \url{https://web.stanford.edu/~aauclert/annual_review_hank.pdf}.
\end{itemize}
Heterogeneity matters to policy analysis; minimal continuous distributional compression is an editorial question.
\begin{enumerate}\raggedright
\item\hypertarget{definitions:source:30007515:1}{} Janet L. Yellen. 2016. \emph{Macroeconomic Research After the Crisis}. Section Heterogeneity.
\url{https://www.federalreserve.gov/newsevents/speech/yellen20161014a.htm}.
\item\hypertarget{definitions:source:30007515:2}{} Adrien Auclert; Matthew Rognlie; Ludwig Straub. 2025. \emph{Fiscal and Monetary Policy with Heterogeneous Agents}. Section5, printed pp.19--21, especially Optimal policy; this does not state a minimal-compression conjecture.
\url{https://web.stanford.edu/~aauclert/annual_review_hank.pdf}.
DOI: \url{https://doi.org/10.1146/annurev-economics-091624-044646}.
\end{enumerate}

\subsection*{Common conventions: Source mathematical conventions}

These conventions apply to each entry unless explicitly replaced.

\textbf{Models and identification.} A primitive model \(m\) specifies all probability laws, feasible choices, information, equilibrium and measurement rules used by its target. The admissible class \(\mathcal M\) is fixed independently of outcomes. If \(P_O(m)\) is its observable probability law and \(\tau(m)\) the target, its identified set at observable law \(P\) is
\[
 \mathcal I_\tau(P)=\{\tau(m):m\in\mathcal M,\ P_O(m)=P\}.
\]
Point identification means that this set is a singleton. An empty set signals incompatibility of data and assumptions. Coordinate bounds need not describe the joint set. Bounds are sharp only if every retained target is attainable or arbitrarily approachable by compatible models. Finite bounds require explicit restrictions on unobserved quantities; unrestricted confounding, hidden wealth, extrapolation or arbitrary equilibrium selection can yield uninformative sets.

\textbf{Causal interventions.} A potential outcome \(Y_i(p)\) is the outcome under a complete policy regime \(p\), including timing, financing, information and market spillovers. The target population and units are fixed before treatment. With interference, use the complete assignment vector or a justified exposure mapping; individual no-interference notation alone cannot identify an economy-wide intervention. A difference of mean potential outcomes does not require the cross-world joint distribution. Conditioning on post-treatment migration, survival, enrollment or employment changes the estimand and needs separate justification.

\textbf{Statistical statements.} An estimator, confidence set, forecast or test specifies a sampling law, model class and loss/coverage criterion. Uniform \((1-\alpha)\)-coverage means \(\inf_{m\in\mathcal M}\Pr_m\{\tau(m)\in C_n\}\geq1-\alpha\), or an explicitly asymptotic version. The whole parameter space trivially has coverage, so informativeness requires a stated width, risk or power objective. A forecasting improvement is relative to a named benchmark, information set and evaluation population.

\textbf{Equilibrium and optimization.} An equilibrium comprises feasible plans, optimality, beliefs where needed, budget constraints and all market/resource clearing. The primitives or a stated benchmark define each correspondence \(\mathcal E(p;m)\). Nonemptiness is assumed or proved, never inferred just from compact policy sets. A selected-equilibrium target uses a declared selection rule; a robust target quantifies over all permitted equilibria. Use suprema or infima unless attainment is established. An \(\varepsilon\)-optimal policy is feasible and within \(\varepsilon\) of the finite optimum value. Report an empty feasible set separately.

\textbf{Welfare and accounting.} Normative utility, interpersonal normalization, social weights, numeraire, discounting and the use of public receipts are primitives. Behavioral utility need not coincide with the welfare criterion. Household welfare includes ownership income and rents once; adding profits again to compensating variation generally double-counts them. A monetary equivalent is defined by an explicit transfer and price convention, with monotonicity and range conditions. Average effects, mechanism contributions, transfers and welfare losses are different targets. Infinite sums require convergence and dynamic feasible plans require terminal or no-Ponzi conditions.

\textbf{Source versions.} A working-paper date, revision date and journal publication year are distinguished. A DOI for a journal version is not automatically the DOI of an earlier manuscript. Locators refer to the stated version and printed pages where specified. A metadata-only check does not verify a claim about a full-text conclusion. Every entry's source subsection narrows the provenance claim accordingly.
% END ECONOMICS STATEMENT
\end{document}
